\section{Related Work}
\label{sec:related}

\paragraph{Graph and Structural-Entropy Social Event Detection.}
Modern offline and incremental SED methods model social posts as heterogeneous message graphs linking shared entities, users, hashtags, and dense semantic embeddings~\cite{cao2021kpgnn,reimers2019sentence}. To avoid brittle supervised edge labels, \texttt{HISEvent}~\cite{cao2024hisevent} introduced greedy hierarchical minimization of 2D Structural Entropy (2D-SE)~\cite{li2016structural} for unsupervised event clustering, later extended into hyperbolic embeddings (\texttt{HyperSED}~\cite{yu2025hypersed}) and relational prompt representations (\texttt{RPLM-SED}~\cite{li2024rplmsed}). Recent LLM-augmented text frameworks inject semantic normalization or active querying into graphs and hierarchies, including \texttt{LSED}~\cite{qiu2025lsed}, \texttt{LHG}~\cite{qiu2026lhg}, \texttt{RagSEDE}~\cite{liu2026ragsede}, and cost-effective pairwise LLM clustering (\texttt{Cequel}~\cite{wang2025cequel}). Our Stage~1 solver inherits exact full-volume greedy 2D-SE agglomerative minimization and local node refinement from this lineage, while our Stage~2 router adapts mesoscopic cluster-pair (\texttt{MERGE\_PAIR}) and minority-pocket (\texttt{SPLIT\_SEAM}, via spherical $k$-means++~\cite{arthur2007kmeanspp}) seam selection to the multimodal VLM setting.

\paragraph{Multimodal SED Benchmarks and Quality Filtering.}
Adding visual features complicates both graph construction and evaluation. On English Twitter, \texttt{NED}~\cite{lin2024ned} pairs $17{,}366$ posts with images across $40$ news events spanning 2011--2021; on Chinese Weibo, \texttt{KQGNN}~\cite{han2026kqgnn} introduces \texttt{WEIBO2022} alongside a key-instance, quality-aware GNN for real-time multimodal SED. Separately, \texttt{CrisisMMD}~\cite{alam2018crisismmd} provides humanitarian informativeness and damage annotations for disaster tweets. We emphasize two important boundary distinctions noted in prior literature reviews: first, \texttt{CrisisMMD}'s \texttt{informative} labels measure emergency-response utility rather than \emph{inter-event discriminability} (generic flood damage can be informative for responders yet non-discriminative between two concurrent hurricanes); second, \texttt{SeInEvent}~\cite{yang2026seinevent} introduces Pointwise Structural Entropy to filter noisy multimodal samples during alternating incremental contrast and rehearsal---rather than discarding test posts at inference---and evaluates open-world \texttt{NED} blocks partitioned chronologically by year. Crucially, our 72-hour chronological audit shows that standard unaligned \texttt{NED} streams yield $B^3\text{-}F_1 = 0.9900$ purely from timestamps, motivating event-relative concurrent evaluation.

\paragraph{Multimodal Interaction Analysis and Vision-Language Models.}
Contrastive language-image pre-training (\texttt{CLIP}~\cite{radford2021clip}) supplies strong zero-shot cross-modal similarity but lacks explicit reasoning when an image matches a topical surface cue from a competing event. Information-theoretic Partial Information Decomposition (\texttt{PID}~\cite{liang2023pid}) formalizes redundancy, modal uniqueness, and synergy with respect to a target variable $Y$. We deliberately refrain from claiming unsupervised neural PID heads---which require target $Y$ identification and synergy estimation that heuristics alone cannot guarantee---and instead adopt \emph{competitor-conditioned causal visual evidence} anchored directly to empirical rescue/mislead counts and conflict-gated \texttt{Qwen/Qwen3-VL-8B-Instruct}~\cite{bai2025qwen3vl} structured adjudication.
